\documentclass[11pt,a4paper]{article}

\usepackage[margin=2.4cm]{geometry}
\usepackage[T1]{fontenc}
\usepackage{lmodern}
\usepackage{microtype}
\usepackage{amsmath,amssymb}
\usepackage{booktabs}
\usepackage{longtable}
\usepackage{array}
\usepackage{xcolor}
\usepackage{listings}
\usepackage{hyperref}

\hypersetup{
  colorlinks=true,
  linkcolor=blue!55!black,
  urlcolor=blue!55!black,
  pdftitle={Numerical Audit and Repair of the 1D3V Electromagnetic PIC Solver}
}

\lstset{
  basicstyle=\ttfamily\small,
  frame=single,
  breaklines=true,
  columns=fullflexible,
  keywordstyle=\color{blue!60!black},
  commentstyle=\color{green!40!black},
  showstringspaces=false
}

\newcommand{\dd}{\mathrm{d}}
\newcommand{\epsz}{\varepsilon_0}
\newcommand{\mucz}{\mu_0}
\newcommand{\code}[1]{\texttt{#1}}

\title{Numerical Audit and Repair of the\\
1D3V Electromagnetic Particle-in-Cell Solver}
\author{Collisionless-Shocks-in-1D-and-2D codebase\\
Technical diagnostic and implementation report}
\date{12 August 2026}

\begin{document}
\maketitle

\begin{abstract}
This report documents a numerical and software audit of a one-dimensional,
three-velocity-component electromagnetic particle-in-cell (PIC) solver.  The
original electrostatic 1D1V solver was stable, whereas the electromagnetic
extension exhibited exponential growth of the transverse components
$E_y,E_z,B_y,$ and $B_z$, eventually contaminating the longitudinal field.
The audit identified a decisive source-free instability in the temporal
composition of the Maxwell update.  Additional independent defects were found
in the relativistic Boris pusher, interpolation of the spatially staggered
magnetic field, current deposition, Gauss-law initialization, macro-particle
normalization, boundary handling, and energy diagnostics.

The repaired periodic 1D3V path uses a stable time-centered Maxwell update,
scalar relativistic Lorentz factors, correct staggered interpolation,
fixed-density macro-particle weights, a longitudinal current satisfying the
discrete continuity equation, an initial periodic Gauss solve, constraint
diagnostics, and relativistic kinetic energy.  Seven regression tests pass.
An end-to-end resolved plasma-oscillation run measured
$\omega_{pe}=0.9967$ for the normalized target $\omega_{pe}=1$, preserved
Gauss's law to $7.36\times10^{-15}$, preserved continuity to
$1.47\times10^{-11}$, and exhibited a maximum relative total-energy drift of
approximately $3.1\times10^{-3}$ over the test interval.  The validated path
is periodic.  Open or absorbing electromagnetic shock boundaries remain
explicitly unsupported pending a separately verified field absorber and
particle injection scheme.
\end{abstract}

\tableofcontents
\newpage

\section{Purpose and scope}

The purpose of this work was to answer the following question using the actual
implementation rather than a generic PIC stability argument:

\begin{quote}
Why did the transverse electromagnetic field components grow exponentially,
and what minimum set of changes is required to distinguish a software defect
from an unstable discretization or under-resolved plasma physics?
\end{quote}

The audit covered the active 1D3V timestep, spatial and temporal staggering,
the Maxwell equations, charge and current deposition, field interpolation,
the relativistic Boris pusher, boundary conditions, energy accounting,
Gauss's law, discrete charge continuity, normalized plasma scales, and minimal
benchmark problems.  The repaired code intentionally validates the periodic
1D3V configuration first.  It does not claim that the existing open-boundary
shock setup is ready for physical production runs.

\section{Code structure}

The principal files are summarized in Table~\ref{tab:files}.

\begin{longtable}{p{0.27\textwidth}p{0.66\textwidth}}
\caption{Main files in the 1D3V implementation.}\label{tab:files}\\
\toprule
File & Responsibility \\
\midrule
\endfirsthead
\toprule
File & Responsibility \\
\midrule
\endhead
\path{solver_1D3V.py} & Initialization and main electromagnetic PIC timestep. \\
\path{grids.py} & Grid allocation, particle-to-grid density deposition, and CIC interpolation indices and weights. \\
\path{maxwell.py} & Current deposition, Gauss initialization and residual, continuity-preserving current construction, and Maxwell updates. \\
\path{newton.py} & Particle drift and relativistic Boris momentum update. \\
\path{particles.py} & Particle state, charge-to-mass ratio, macro-particle weight, velocity--momentum conversion, and kinetic energy. \\
\path{boundary_conditions.py} & Particle boundary operations. \\
\path{time_constraint.py} & Electromagnetic CFL and plasma-frequency timestep estimates. \\
\path{results.py} & Field, particle, energy, charge, continuity, and Gauss diagnostics. \\
\path{runs/periodic_1d3v.py} & Resolved periodic end-to-end example. \\
\path{tests/test_1d3v_core.py} & Seven numerical regression tests. \\
\bottomrule
\end{longtable}

\section{Governing one-dimensional Maxwell system}

In 1D3V all fields have three components, but all spatial derivatives are in
$x$.  The continuous equations used by the code are
\begin{align}
  \frac{\partial E_x}{\partial t} &= -\frac{J_x}{\epsz}, \\
  \frac{\partial E_y}{\partial t} &= -\frac{J_y}{\epsz}
      -c^2\frac{\partial B_z}{\partial x}, \\
  \frac{\partial E_z}{\partial t} &= -\frac{J_z}{\epsz}
      +c^2\frac{\partial B_y}{\partial x}, \\
  \frac{\partial B_x}{\partial t} &= 0, \\
  \frac{\partial B_y}{\partial t} &= \frac{\partial E_z}{\partial x}, \\
  \frac{\partial B_z}{\partial t} &= -\frac{\partial E_y}{\partial x}.
\end{align}
The signs in the active periodic implementation were correct.  The transverse
instability arose from the temporal update order, not from a sign reversal in
Faraday's or Amp\`ere--Maxwell's law.

The spatial Yee pairing places transverse electric samples at integer grid
locations and transverse magnetic samples half a cell away.  The corresponding
periodic differences are
\begin{align}
  (D_- B)_i &= \frac{B_i-B_{i-1}}{\Delta x}, \\
  (D_+ E)_i &= \frac{E_{i+1}-E_i}{\Delta x}.
\end{align}
These operators are the correct adjoint pair on a periodic grid.

\section{Original timestep and staggering}

At the beginning of an original loop iteration the intended particle state was
$x_p^n$ and $v_p^{n+1/2}$.  The code then performed the following operations:
\begin{enumerate}
  \item deposit a current at $x_p^n$;
  \item drift particles to $x_p^{n+1}$;
  \item apply particle boundary conditions;
  \item deposit density and another current at $x_p^{n+1}$;
  \item average the two endpoint currents as an approximation to
        $J^{n+1/2}$;
  \item advance all components of $E$ and $B$;
  \item Boris-push momentum from the old to the new particle half-step.
\end{enumerate}

The endpoint average is formally centered for smooth trajectories that remain
inside one cell, but it does not enforce the discrete continuity equation when
a particle changes its CIC support.  Consequently it cannot guarantee that an
initially satisfied Gauss constraint remains satisfied under the longitudinal
Amp\`ere update.

\section{Confirmed source-free Maxwell instability}
\label{sec:maxwell-instability}

\subsection{Original temporal composition}

The original routine divided a timestep into two nominal half-steps but used
old and partially updated fields in a non-symplectic order.  For one transverse
polarization, write the semi-discrete vacuum equations as
\begin{equation}
  \dot E=aB,\qquad \dot B=bE,
\end{equation}
where $a=-c^2D_-$ and $b=-D_+$.  For a Fourier mode,
$ab=-\Omega_k^2$, with
\begin{equation}
  \Omega_k = \frac{2c}{\Delta x}
  \sin\left(\frac{k\Delta x}{2}\right).
\end{equation}
Let $h=\Delta t/2$.  The original source-free sequence was equivalent to
\begin{align}
  E^* &= E^n+h aB^n,\\
  B^{n+1/2} &= B^n+h bE^n,\\
  B^{n+1} &= B^{n+1/2}+h bE^*,\\
  E^{n+1} &= E^*+h aB^{n+1/2}.
\end{align}
Therefore
\begin{equation}
\begin{bmatrix}E^{n+1}\\B^{n+1}\end{bmatrix}
=
\begin{bmatrix}
1+h^2ab & 2ha\\
2hb & 1+h^2ba
\end{bmatrix}
\begin{bmatrix}E^n\\B^n\end{bmatrix}.
\end{equation}
Defining $s=h^2\Omega_k^2>0$, the eigenvalue magnitude is
\begin{equation}
  |\lambda_k|=1+s>1\qquad(k\ne0).
\end{equation}
Thus every nonzero vacuum mode grows.  Reducing $\Delta t$ reduces the growth
per step but does not change its sign.  PIC shot noise supplies broadband
transverse seeds, and the grid-scale modes grow fastest.  Disabling magnetic
evolution removes the $E\leftrightarrow B$ feedback loop, explaining the
observed suppression of transverse growth when $B$ was held fixed.

\subsection{Direct reproduction}

For $N_x=200$, $\Delta x=0.005$, $\Delta t=0.00025$,
$c\Delta t/\Delta x=0.05$, zero current, and a mode-50 transverse electric
perturbation, the original field energy evolved as shown in
Table~\ref{tab:old-energy}.

\begin{table}[h]
\centering
\caption{Source-free energy growth under the original update.}
\label{tab:old-energy}
\begin{tabular}{rr}
\toprule
Step & Field energy \\
\midrule
0 & 0.250000 \\
1 & 0.250625 \\
100 & 0.320956 \\
500 & 0.871905 \\
1000 & 3.040872 \\
\bottomrule
\end{tabular}
\end{table}

This test contains no particles, no current deposition, no Debye length, and
no physical plasma instability.  It proves that the Maxwell update alone was
sufficient to generate the reported behavior.

\section{Maxwell repair}

The repaired routine in \path{maxwell.py} uses a centered
$B/2\rightarrow E\rightarrow B/2$ composition:
\begin{align}
  B^{n+1/2} &= B^n + \frac{\Delta t}{2}F(E^n),\\
  E^{n+1} &= E^n + \Delta t\left[G(B^{n+1/2})
               -\frac{J^{n+1/2}}{\epsz}\right],\\
  B^{n+1} &= B^{n+1/2}+\frac{\Delta t}{2}F(E^{n+1}),
\end{align}
where $F$ represents Faraday's curl and $G$ represents the magnetic part of
Amp\`ere--Maxwell's law.  This is the velocity-Verlet form of the spatially
staggered wave system.  It has bounded source-free solutions under the 1D CFL
condition
\begin{equation}
  \frac{c\Delta t}{\Delta x}\leq 1.
\end{equation}

Using the same mode-50 test for 2000 steps, the repaired field energy remained
between $0.2496875001$ and $0.2499999999$.  The small oscillation is bounded
time-discretization error rather than exponential energy injection.

\section{Relativistic Boris pusher defects and repair}

\subsection{Original defects}

The particle momentum is $\mathbf{u}=\gamma\mathbf{v}$.  The original code
computed
\begin{equation}
  \gamma^{-1}_{p,j}=\frac{1}{\sqrt{1+u_{p,j}^2}},
\end{equation}
separately for each Cartesian component.  Relativistic gamma is instead a
single scalar for each particle, determined by the norm of its momentum:
\begin{equation}
  \gamma_p=\sqrt{1+\frac{|\mathbf{u}_p|^2}{c^2}}.
\end{equation}
The old implementation also used the gamma computed before magnetic rotation
and before the second electric impulse to convert the final momentum back to
velocity.  The initialization method \code{v\_to\_u} contained the analogous
componentwise error.

Because the algebraic Boris rotation preserves $|\mathbf{u}|$ for a fixed
rotation vector, the internal momentum norm could look correct while the
reported velocity, trajectory, current, gyrofrequency, and kinetic energy were
wrong.  In a uniform-$B$ test with initial speed $0.6c$, the old code produced
speeds between $0.6000c$ and $0.6626c$ even though $E=0$.

\subsection{Corrected update}

The repaired routine performs
\begin{align}
  \mathbf{u}^{-} &= \mathbf{u}^{n}
       +\frac{q\Delta t}{2m}\mathbf{E},\\
  \gamma^- &= \sqrt{1+|\mathbf{u}^{-}|^2/c^2},\\
  \mathbf{t} &= \frac{q\Delta t}{2m\gamma^-}\mathbf{B},\\
  \mathbf{s} &= \frac{2\mathbf{t}}{1+|\mathbf{t}|^2},\\
  \mathbf{u}' &= \mathbf{u}^{-}+\mathbf{u}^{-}\times\mathbf{t},\\
  \mathbf{u}^{+} &= \mathbf{u}^{-}+\mathbf{u}'\times\mathbf{s},\\
  \mathbf{u}^{n+1} &= \mathbf{u}^{+}
       +\frac{q\Delta t}{2m}\mathbf{E},\\
  \gamma^{n+1} &= \sqrt{1+|\mathbf{u}^{n+1}|^2/c^2},\\
  \mathbf{v}^{n+1} &= \mathbf{u}^{n+1}/\gamma^{n+1}.
\end{align}
Electron and ion rotation directions follow automatically from their signed
charge-to-mass ratios.  The repaired uniform-$B$ regression conserves both
$|\mathbf{u}|$ and speed $0.6c$ to numerical precision over 1000 steps.

\section{Staggered magnetic interpolation}

If $B_j$ is located at $(j+1/2)\Delta x$, a particle at $x_p$ is bracketed by
\begin{equation}
  j=\left\lfloor\frac{x_p}{\Delta x}-\frac12\right\rfloor,
\end{equation}
and its correct linear weight is
\begin{equation}
  w_B=\frac{x_p}{\Delta x}-\left(j+\frac12\right).
\end{equation}
The original implementation omitted the $1/2$ in the weight and used
$x_p/\Delta x-j$.  Its weights therefore lay in $[0.5,1.5)$ instead of
$[0,1)$, amounting to a half-cell sampling shift in the interior and
extrapolation near periodic wrapping.  This error altered magnetic forces and
could exchange spurious energy with particles.

The corrected expressions are implemented for both species in
\path{grids.py}.  A regression using a linear magnetic field recovers the
particle coordinate exactly.

\section{Macro-particle normalization}

Originally each simulation particle deposited its literal charge without a
macro-particle weight or division by $\Delta x$.  Consequently the mean density
and plasma frequency changed when the number of particles per cell changed.
This made particle-count convergence physically ill-defined.

The repaired periodic solver introduces a normalized continuum density $n_0$
in \path{parameters.py}.  For a one-dimensional unit-area domain of length
$L$, each species uses
\begin{equation}
  w_s=\frac{n_0L}{N_s}.
\end{equation}
Number density is deposited as
\begin{equation}
  n_{s,i}=\frac{w_s}{\Delta x}
      \sum_{p\in s}S_i(x_p),
\end{equation}
and current density carries the corresponding factor
$q_sw_s/\Delta x$.  Particle acceleration continues to use the physical
$q_s/m_s$, so changing $N_s$ now changes statistical noise without changing
the intended continuum density.  A regression with 32 and 128 particles
verifies that both cases integrate to the same density $n_0=1$.

For the normalized constants $|q_e|=m_e=\epsz=1$,
\begin{equation}
  \omega_{pe}=\sqrt{n_0}.
\end{equation}
The timestep estimator now includes $n_0$ in its plasma-frequency estimate.

\section{Charge continuity and Gauss's law}

\subsection{Problem in the endpoint-average current}

For consistent longitudinal electromagnetic evolution, the deposited charge
and current must satisfy the same discrete identity used by Amp\`ere's law:
\begin{equation}
  \frac{\rho_i^{n+1}-\rho_i^n}{\Delta t}
  +\frac{J_{x,i}^{n+1/2}-J_{x,i-1}^{n+1/2}}{\Delta x}=0.
  \label{eq:continuity}
\end{equation}
An average of CIC currents deposited at trajectory endpoints does not satisfy
Eq.~\eqref{eq:continuity} exactly when a particle changes its supporting grid
nodes.  Thus it permits longitudinal Gauss error to be generated at every
step.

\subsection{Periodic 1D continuity-preserving construction}

In one periodic dimension, Eq.~\eqref{eq:continuity} determines $J_x$ up to a
spatially uniform constant.  Let
$\Delta\rho_i=\rho_i^{n+1}-\rho_i^n$.  Starting from a provisional
$J_{x,0}=0$, the implementation constructs
\begin{equation}
  J_{x,i}=-\frac{\Delta x}{\Delta t}
      \sum_{m=1}^{i}\Delta\rho_m,
  \qquad i=1,\ldots,N_x-1.
\end{equation}
Neutrality makes the remaining periodic equation at $i=0$ consistent.  The
undetermined uniform part is then chosen from the total particle trajectory
flux,
\begin{equation}
  \langle J_x\rangle=rac{1}{L\Delta t}
       \sum_p q_pw_p\Delta x_p,
\end{equation}
where $\Delta x_p$ is the shortest periodic displacement.  This procedure is
specific to 1D periodic geometry; it should not be presented as a replacement
for a general multidimensional Esirkepov-type deposition algorithm.

The transverse currents $J_y$ and $J_z$ do not enter the 1D charge-continuity
divergence.  They are deposited by CIC at the periodic midpoint of each
particle trajectory, using velocities at the particle half-step.

\subsection{Initial and preserved Gauss constraint}

The 1D3V grid now stores
\begin{equation}
  \rho_i=q_en_{e,i}+q_in_{i,i}.
\end{equation}
For a periodic neutral domain, the initial longitudinal electric field is
constructed to satisfy
\begin{equation}
  \frac{E_{x,i}-E_{x,i-1}}{\Delta x}=\frac{\rho_i}{\epsz}
\end{equation}
to roundoff.  Its spatially uniform, undetermined component is set to zero.
If the periodic domain has nonzero net charge, initialization raises an error
because no periodic solution of Gauss's law exists.

Under the longitudinal update
\begin{equation}
 E_x^{n+1}=E_x^n-\frac{\Delta t}{\epsz}J_x^{n+1/2},
\end{equation}
Eq.~\eqref{eq:continuity} causes discrete Gauss's law to be preserved
algebraically.  The code records both $L_2$ and $L_\infty$ norms of the
continuity and Gauss residuals.

\section{Energy accounting and diagnostics}

The original particle diagnostic used the classical expression
$\tfrac12m|\mathbf{v}|^2$ even for $v=0.9c$.  The periodic electromagnetic path
now uses the weighted relativistic kinetic energy
\begin{equation}
  K_s=\sum_{p\in s}w_sm_sc^2(\gamma_p-1).
\end{equation}
The field energy remains
\begin{equation}
  U_F=\frac{\Delta x}{2}\sum_i
  \left(\epsz|\mathbf{E}_i|^2+rac{|\mathbf{B}_i|^2}{\mucz}\right).
\end{equation}

The results object additionally stores
\begin{itemize}
  \item charge density $\rho$;
  \item Gauss residual $L_2$ and $L_\infty$ norms;
  \item continuity residual $L_2$ and $L_\infty$ norms;
  \item longitudinal electric energy;
  \item transverse electric energy;
  \item magnetic energy.
\end{itemize}
This separation makes it possible to identify whether energy first appears in
the longitudinal particle coupling or in the transverse Maxwell subsystem.

\section{Boundary conditions}

The old driver accepted open or absorbing particle boundaries but always
called a periodic field solver.  This silently coupled nonperiodic particles to
periodically wrapping fields.  In addition, the unused nonperiodic field path
contained endpoint indexing and zero-locking limiter defects.

The repaired \path{solver_1D3V.py} explicitly accepts only
\code{BoundaryCondition.Periodic}.  Open and absorbing choices raise
\code{NotImplementedError}.  This is deliberate: failing clearly is safer than
returning a plausible-looking but mathematically inconsistent shock result.
The periodic particle wrapper was also changed so it no longer subtracts a
small displacement from every particle at every step.

A production open shock calculation still requires
\begin{enumerate}
  \item a verified outgoing-wave or perfectly matched electromagnetic layer;
  \item charge-consistent particle removal and injection;
  \item synchronized creation and deletion of $x$, $v$, and $u$;
  \item boundary contributions to the energy and charge budgets;
  \item reflection and transmission benchmark tests.
\end{enumerate}

\section{Resolution and normalized parameters}

For normalized $n_0=1$, $|q_e|=m_e=\epsz=c=1$,
\begin{equation}
  \omega_{pe}=1,\qquad d_e=\frac{c}{\omega_{pe}}=1,
  \qquad \lambda_D\simeq\frac{v_{te}}{\omega_{pe}}=v_{te}.
\end{equation}
The historical cold setup used $v_{te}=10^{-8}$ and $\Delta x=0.005$, giving
\begin{equation}
  \frac{\lambda_D}{\Delta x}=2\times10^{-6}.
\end{equation}
This is severely under-resolved.  It can produce finite-grid heating and alias
instabilities, but it was not the cause of the proven source-free Maxwell
growth in Section~\ref{sec:maxwell-instability}.

The resolved example uses $N_x=100$, $\Delta x=0.01$, and $v_{te}=0.01$, so
$\lambda_D/\Delta x\simeq1$.  Its timestep is $\Delta t=5\times10^{-4}$,
giving
\begin{equation}
  \frac{c\Delta t}{\Delta x}=0.05,
  \qquad \omega_{pe}\Delta t=5\times10^{-4}.
\end{equation}
Both conditions are comfortably inside their explicit stability limits.

Using $v_{te}=10^{-8}$ with a unit-length explicit grid would require of order
$10^8$ cells merely to reach $\Delta x\lesssim\lambda_D$, which is not a
practical configuration for this Python PIC code.  A cold physical problem
must therefore use a different normalization, a reduced scale separation, or
an implicit/semi-implicit method rather than simply restoring the old thermal
speed.

\section{Verification suite}

The regression suite is located in \path{tests/test_1d3v_core.py}.  Its current
coverage is summarized in Table~\ref{tab:tests}.

\begin{longtable}{p{0.38\textwidth}p{0.53\textwidth}}
\caption{Implemented numerical regression tests.}\label{tab:tests}\\
\toprule
Test & Acceptance criterion \\
\midrule
\endfirsthead
\toprule
Test & Acceptance criterion \\
\midrule
\endhead
Vacuum transverse mode & Source-free field energy remains bounded over 2000 steps. \\
Uniform magnetic field & Relativistic Boris rotation preserves momentum norm and speed. \\
Staggered magnetic CIC & A linear staggered field interpolates exactly at the particle position. \\
Velocity-to-momentum conversion & Gamma is computed from total speed rather than per component. \\
Macro-particle normalization & Changing particle count does not change integrated or mean density. \\
Continuity and Gauss preservation & A periodic cell-crossing trajectory satisfies both constraints to roundoff. \\
Electron plasma oscillation & Measured normalized frequency agrees with $\omega_{pe}=1$ within three percent and does not create transverse fields. \\
\bottomrule
\end{longtable}

The plasma-oscillation test measured
\begin{equation}
  \omega_{pe,\mathrm{measured}}=0.9967,
\end{equation}
which is consistent with the normalized target and the finite grid and shape
factor.

\section{End-to-end periodic example}

The executable example \path{runs/periodic_1d3v.py} exercises the public
\code{solver\_1D3V.simulate} entry point with 4000 particles, 100 cells,
$n_0=1$, $v_{te}=0.01$, periodic boundaries, and three velocity components.
It ran to $t=2.0005$ in 4001 steps and reported
\begin{align}
  \max\|G\|_\infty &= 7.36\times10^{-15},\\
  \max\|C\|_\infty &= 1.47\times10^{-11},\\
  \max\left|\frac{E_{\mathrm{tot}}(t)-E_{\mathrm{tot}}(0)}
  {E_{\mathrm{tot}}(0)}\right| &= 3.09\times10^{-3},\\
  \max U_{\mathrm{transverse}} &\simeq 1.05\times10^{-7},
\end{align}
where $G$ is the discrete Gauss residual and $C$ is the discrete continuity
residual.  These results establish that the repaired periodic code completes a
coupled particle--field run without the old exponential transverse failure.
The energy result is a benchmark value, not a proof of exact energy
conservation for arbitrary nonlinear distributions.

\section{Current supported status}

The following statement is scientifically defensible:

\begin{quote}
The periodic 1D3V electromagnetic PIC path is numerically functional and has
passed isolated vacuum, pusher, interpolation, normalization, constraint, and
plasma-frequency benchmarks.  The previously observed source-free transverse
exponential instability has been removed.
\end{quote}

The following stronger statement is not yet justified:

\begin{quote}
The code is ready for production open-boundary collisionless shock studies at
the historical cold-plasma parameters.
\end{quote}

The reasons are that electromagnetic open/absorbing boundaries have not been
implemented and validated, and the historical Debye length is unresolved by
many orders of magnitude.

\section{Reproduction instructions}

From the repository root, run the regression suite with
\begin{lstlisting}[language=bash]
python -m unittest discover -s tests -v
\end{lstlisting}
and the resolved periodic example with
\begin{lstlisting}[language=bash]
python runs/periodic_1d3v.py
\end{lstlisting}

The example disables result-file writing by default while still returning all
diagnostics in memory.  Normal simulations retain the existing behavior of
writing results unless \code{write\_results=False} is passed to
\code{simulate}.

\section{Recommended next work}

Before constructing a shock, the following work should be performed in order:
\begin{enumerate}
  \item add longer energy-convergence tests over multiple plasma periods;
  \item add cold-plasma electromagnetic dispersion and two-stream growth-rate
        benchmarks;
  \item vary particles per cell at fixed $n_0$ and verify noise convergence;
  \item vary $\Delta x$ and $\Delta t$ independently while keeping
        $\lambda_D/\Delta x$ controlled;
  \item implement an electromagnetic absorbing boundary and measure its
        reflection coefficient using incident vacuum waves;
  \item design charge-consistent particle inflow and outflow and verify global
        charge and energy budgets;
  \item only then run a shock and separate physical growth rates from numerical
        noise through resolution convergence.
\end{enumerate}

\appendix
\section{Summary of modified files}

\begin{description}
  \item[\path{maxwell.py}] Stable Maxwell composition, periodic Gauss solve,
    Gauss residual, normalized current, midpoint transverse current, and
    continuity-preserving longitudinal current.
  \item[\path{newton.py}] Correct scalar relativistic gamma and final velocity
    conversion in the Boris pusher.
  \item[\path{grids.py}] Correct half-cell magnetic weights, physical number
    density, charge density, and constraint arrays for \code{Grid1D3V}.
  \item[\path{particles.py}] Macro-particle weight, total-speed momentum
    initialization, speed validation, and relativistic kinetic energy.
  \item[\path{solver_1D3V.py}] Fixed-density periodic PIC cycle, constraint
    monitoring, relativistic energy, and explicit boundary support guard.
  \item[\path{parameters.py}] Normalized continuum density $n_0$.
  \item[\path{time_constraint.py}] Density-aware plasma-frequency estimate.
  \item[\path{results.py}] Charge, constraints, and component field energies.
  \item[\path{boundary_conditions.py}] Periodic wrapping without an artificial
    per-step particle displacement.
  \item[\path{tests/test_1d3v_core.py}] Seven numerical regression tests.
  \item[\path{runs/periodic_1d3v.py}] Resolved end-to-end example.
  \item[\path{README.md}] Documented validated scope and commands.
\end{description}

\section{Branch recommendation}

The repairs should be maintained on a separate branch until reviewed and
compared against the original implementation.  A suitable branch and commit
sequence is
\begin{lstlisting}[language=bash]
git switch -c codex/fix-1d3v-pic-stability
git add README.md boundary_conditions.py grids.py maxwell.py newton.py \
  parameters.py particles.py results.py solver_1D3V.py time_constraint.py \
  runs/periodic_1d3v.py tests docs/1D3V_numerical_audit_and_repairs.tex
git commit -m "Fix and validate periodic 1D3V PIC solver"
git push -u origin codex/fix-1d3v-pic-stability
\end{lstlisting}

\end{document}
